\section{Introdução}

A calibração de uma câmera estima os parâmetros que relacionam pontos de uma
cena às suas posições na imagem. A matriz intrínseca reúne as distâncias focais
em pixels e o ponto principal; os coeficientes de distorção representam desvios
em relação ao modelo de projeção ideal. Esses parâmetros permitem corrigir
deformações geométricas e projetar pontos tridimensionais em coordenadas de
imagem, conforme os procedimentos do OpenCV~\cite{opencv-tutorial,opencv-calib3d}.

Neste trabalho, foram calibradas duas câmeras de um smartphone. Os parâmetros
obtidos foram aplicados à correção de fotografias de teste e à projeção de
pontos de um tabuleiro. O principal experimento quantitativo utiliza novas
fotografias, não incluídas na calibração, e compara posições observadas e
previstas de quatro cantos reservados da estimação da pose em cada imagem.
